\documentclass{ctexart}
\newcommand{\keywords}[1]{#1}
\begin{document}
\begin{abstract}
本文针对几何调整问题建立约束模型，给出可复算结果并以残差检验。
\end{abstract}
\keywords{几何；优化；误差}
\section{问题重述}
在给定结构尺寸和行程限制下确定节点调整量。
\section{问题分析}
先建立几何基线，再把结构约束写入优化模型。
\section{模型假设}
忽略测量噪声的高频分量，并在后文检验其影响。
\section{符号说明}
$x_i$ 为节点位移，单位为米。
\section{模型建立与求解}
以最大几何误差最小为目标，显式保留长度和行程约束，采用数值求解。
\section{结果分析}
结果表明最大误差相对基线下降，同时所有约束残差均满足阈值。该变化来自结构兼容条件约束下的联合调整，因此数值改善对应可执行候选。回代后最大长度残差和最大行程残差分别低于题面上界，并且边界节点没有越界，说明结果在当前离散精度和给定参数范围内可行。
\section{模型检验}
回代复算的 RMSE 为给定数值，并与基准结果对照；网格加密后变化低于阈值。
\section{灵敏度分析}
对长度阈值扰动并重新求解，记录目标和活跃约束变化。
\section{模型评价与改进}
模型保留结构约束，但参数误差仍需更多观测校准。
\begin{thebibliography}{1}
\bibitem{ref1} 某作者. 数值方法教材. 2020.
\end{thebibliography}
\appendix
\section{附录：代码与数据}
列出入口、环境和复算命令。
\end{document}
